\chapter{球函数}

\problem{16.1}{试证明 Legendre 多项式的生成函数
\[
  f(r) = \frac{1}{\sqrt{1+r^2-2rx}}
\]
在 $r=0$ 点的 Taylor 展开式为
\[
  f(r) = \sum_{l=0}^{\infty} \mathrm{P}_l(x)r^l
\]
其中 $r=|\boldsymbol{r}|<1,\,x=\cos\theta,\,|x|\leqslant 1$。}
\answerarea

\problem{16.2}{试证明 Legendre 多项式的递推公式：
\begin{enumerate}
  \item $(2l+1)\mathrm{P}_l(x)=\mathrm{P}_{l+1}'(x)-\mathrm{P}_{l-1}'(x)$
  \item $(l+1)\mathrm{P}_l(x)=\mathrm{P}_{l+1}'(x)-x\mathrm{P}_l'(x)$
  \item $l\mathrm{P}_l(x)=x\mathrm{P}_l'(x)-\mathrm{P}_{l-1}'(x)$
  \item $(x^2-1)\mathrm{P}_l'(x)=lx\mathrm{P}_l(x)-l\mathrm{P}_{l-1}(x)$
\end{enumerate}}
\answerarea
\clearpage

\problem{16.3}{试证明：
\begin{enumerate}
  \item $\mathrm{P}_l(1)=1$
  \item $\mathrm{P}_l(-1)=(-1)^l$
  \item $\mathrm{P}_l(0)=\begin{cases} 0 & l=2n+1,\ n=0,1,2,\cdots \\[4pt] (-1)^n\dfrac{(2n-1)!!}{2n!!} & l=2n,\ n=0,1,2,\cdots \end{cases}$
  \item $\mathrm{P}_l'(1)=\dfrac{l(l+1)}{2}$
  \item $\mathrm{P}_l'(-1)=(-1)^{l-1}\dfrac{l(l+1)}{2}$
  \item $\mathrm{P}_l'(0)=\begin{cases} (-1)^n\dfrac{(2n+1)!}{(2n)!!}, & l=2n+1,\ n=0,1,2,\cdots \\[6pt] 0, & l=2n,\ n=0,1,2,\cdots \end{cases}$
\end{enumerate}}
\answerarea

\problem{16.4}{计算下列积分：
\begin{enumerate}
  \item $\displaystyle\int_{-1}^{1} \mathrm{P}_l(x)\,\mathrm{d}x$ \qquad
  \item $\displaystyle\int_{-1}^{1} x^2\mathrm{P}_l(x)\mathrm{P}_{l+2}(x)\,\mathrm{d}x$
  \item $\displaystyle\int_{-1}^{1} (1-x^2)[\mathrm{P}_l'(x)]^2\,\mathrm{d}x$ \qquad
  \item $\displaystyle\int_{-1}^{1} x\mathrm{P}_l(x)\mathrm{P}_{l+1}(x)\,\mathrm{d}x$
  \item $\displaystyle\int_{0}^{1} \mathrm{P}_l(x)\,\mathrm{d}x$ \qquad
  \item $\displaystyle\int_{0}^{1} \mathrm{P}_k(x)\mathrm{P}_l(x)\,\mathrm{d}x$
  \item $\displaystyle\int_{-1}^{1} x^k\mathrm{P}_l(x)\,\mathrm{d}x$，分别讨论 $k\geqslant l$ 和 $k<l$ 的情况。
\end{enumerate}}
\answerarea
\clearpage

\problem{16.5}{在以 Legendre 多项式为基 $\{\mathrm{P}_l(x),\,l=0,1,2,\cdots\}$ 的函数空间中，将下列函数展开。
\begin{enumerate}
  \item $f(x)=|x|$， \qquad
  \item $f(x)=1+x+x^2+x^3$
\end{enumerate}}
\answerarea

\problem{16.6}{在点电荷 $4\pi\varepsilon_0 q$ 的电场中放置一半径为 $a$ 的导体球壳，球心距点电荷的距离为 $b$，且 $b>a$，求这体系所形成的电场。}
\answerarea
\clearpage

\problem{16.7}{半径为 $a$ 的匀质半球，其半球面的温度保持恒定为 $T_0$，其底面温度保持零度，求球内温度的稳定分布。}
\answerarea

\problem{16.8}{求以下两种均匀带电体在空间形成的稳定电场分布。
\begin{enumerate}
  \item 均匀带电的薄圆盘，圆盘的半径为 $a$，总电荷量为 $4\pi\varepsilon_0 q$。
  \item 均匀带电的细圆环，环的半径为 $a$，总电荷量为 $4\pi\varepsilon_0 q$。
\end{enumerate}}
\answerarea
\clearpage

\problem{16.9}{求解半径为 $a$ 的空心球壳内的定解问题
\[
  \left\{\begin{aligned}
    &\nabla^2 u = 0,\quad r<a \\
    &u\big|_{r=a} = \cos^2\theta,\quad 0\leqslant\theta\leqslant\pi
  \end{aligned}\right.
\]}
\answerarea

\problem{16.10}{试证明
\[
  \mathrm{P}_l^{-m}(x) = (-1)^m \frac{(l-m)!}{(l+m)!}\mathrm{P}_l^m(x),\quad m>0
\]}
\answerarea
\clearpage

\problem{16.11}{试证明
\[
  \mathrm{P}_l^m(0) = \begin{cases}
    (-1)^{\frac{1}{2}(l-m)}\dfrac{(l+m-1)!!}{(l-m)!!}, & \text{当 }l+m\text{ 为偶数} \\[8pt]
    0, & \text{当 }l-m\text{ 为奇数}
  \end{cases}
\]}
\answerarea

\problem{16.12}{试证明
\[
  \mathrm{P}_l^l(\cos\theta) = (2l-1)!!\sin^l\theta,\quad l=0,1,2,\cdots
\]}
\answerarea
\clearpage

\problem{16.13}{在以球函数为基 $\left\{\mathrm{Y}_{l,m}(\theta,\varphi),\,\begin{aligned} &l=0,1,2,\cdots \\ &m=0,\pm 1,\cdots,\pm l \end{aligned}\right\}$ 的函数空间中，将下列函数展开
\begin{enumerate}
  \item $\sin^2\theta\cos^2\varphi$ \qquad
  \item $(1+3\cos\theta)\sin\theta\cos\varphi$
  \item $\dfrac{1}{r^2}(x^2+2z^2+3xy+4xy)$
\end{enumerate}
其中 $r^2=x^2+y^2+z^2,\,0\leqslant\theta\leqslant\pi,\,0\leqslant\varphi\leqslant 2\pi$。}
\answerarea

\problem{16.14}{在半径为 $a$ 的球壳内外，分别求解下面定解问题的球外区域和球内区域的解
\[
  \left\{\begin{aligned}
    &\nabla^2 u = 0 \\
    &u\big|_{r=a} = f(\theta,\varphi) \\
    &u\big|_{r=0}<\infty,\quad u\big|_{r\to\infty}<\infty
  \end{aligned}\right.
\]
并分别对以下两种 $f(\theta,\varphi)$ 的形式
\begin{enumerate}
  \item $f(\theta,\varphi)=\cos\theta$，
  \item $f(\theta,\varphi)=\mathrm{P}_l^1(\cos\theta)\cos\varphi$，给出解的具体形式。
\end{enumerate}}
\answerarea
\clearpage
